\section{Fixed bounded formulas, truth bits, and finite strategies}
\label{sec:bounded}

Bounded quantifiers are mathematically finite, but they are still quantifiers.
They therefore cannot appear directly in an induction matrix restricted to
quantifier-free formulas.  DARC removes them before induction by compiling each
fixed bounded formula to a primitive-recursive truth bit.  When a proof needs
witness information rather than truth alone, the same syntax generates finite
strategy certificates.

\subsection{Reflected atoms}

An atomic interface consists of:
\begin{enumerate}
  \item a formula \(A(\vec x)\);
  \item a primitive-recursive bit \(\chi_A(\vec x)\in\{0,1\}\);
  \item ordinary derivations of bit range and reflection:
\[
  \vdash \chi_A(\vec x)=0\vee\chi_A(\vec x)=1,
  \qquad
  \vdash A(\vec x)\leftrightarrow \chi_A(\vec x)=1.
\]
\end{enumerate}
Equality of terms, order, divisibility with a fixed convention, finite code
membership, and transition relations are typical reflected atoms.  The atom
interface is theorem-specific: DARC does not assert that every arbitrary
formula has a characteristic bit.

\subsection{Boolean operations}

Assume \(a,b\in\{0,1\}\).  Primitive-recursive arithmetic can realize:
\begin{align*}
  a\wedge b &: =ab,\\
  a\vee b &: =\operatorname{sg}(a+b),\\
  \neg a &: =1\monusop a,\\
  a\Rightarrow b &: =\operatorname{sg}((1\monusop a)+b),
\end{align*}
where \(\monusop\) is truncated subtraction and
\(\operatorname{sg}(z)=0\) iff \(z=0\), otherwise \(1\).  On bits, the truth
tables are proved by the decidable zero/successor decomposition available in
the arithmetic library.

The implication clause deserves comment.  Intuitionistically,
\((P\to Q)\leftrightarrow(\neg P\vee Q)\) is not valid for arbitrary \(P\).
It is valid when \(P\) is decidable.  The reflected bit gives precisely that
decidability.  Thus the compiler may support implication in bounded formulas
without importing unrestricted excluded middle.

\subsection{Finite folds for bounded quantifiers}

Let \(b(i,\vec x)\in\{0,1\}\).  Define
\begin{align*}
  \operatorname{All}(0,\vec x)&=1,
  &\operatorname{All}(S n,\vec x)&=
      \operatorname{All}(n,\vec x)b(n,\vec x),\\
  \operatorname{Any}(0,\vec x)&=0,
  &\operatorname{Any}(S n,\vec x)&=
      \operatorname{sg}(\operatorname{Any}(n,\vec x)+b(n,\vec x)).
\end{align*}
Both are primitive recursive.  They represent bounded universal and
existential quantification.

\begin{lemma}[Finite universal reflection]
\label{lem:all-reflect}
If \(b(i,\vec x)=1\leftrightarrow B(i,\vec x)\), then
\[
 \operatorname{All}(N,\vec x)=1
 \quad\leftrightarrow\quad
 \forall i<N\;B(i,\vec x).
\]
The displayed bounded universal is an abbreviation for its ordinary
first-order formulation, not an induction matrix.
\end{lemma}

\begin{proof}
Prove the finite-fold invariant by quantifier-free induction on \(N\): the
product is one exactly when every factor so far is one.  The base is the empty
product.  At the successor step, use the bit range and multiplication truth
table to split
\(\operatorname{All}(S N)=1\) into
\(\operatorname{All}(N)=1\) and \(b(N)=1\).  Translate the latter by atomic
reflection.  The final bounded-quantifier statement can then be assembled by
ordinary quantifier rules or retained as the reflected bit, depending on the
consumer.
\end{proof}

\begin{lemma}[Finite existential reflection]
\label{lem:any-reflect}
If \(b(i,\vec x)=1\leftrightarrow B(i,\vec x)\), then
\[
 \operatorname{Any}(N,\vec x)=1
 \quad\leftrightarrow\quad
 \exists i<N\;B(i,\vec x).
\]
\end{lemma}

\begin{proof}
The forward direction can use bounded least search on the predicate
\(1\monusop b(i,\vec x)=0\).  The recurrence for \(\operatorname{Any}\)
proves that if the final bit is one, some factor bit is one; search chooses the
least such index.  The reverse direction proves by induction that a true bit at
an index below \(N\) makes the finite disjunction bit one.  All induction
matrices are equations and inequalities between primitive-recursive terms.
\end{proof}

\subsection{Structural reflection theorem}

Fix a bounded-formula syntax generated from reflected atoms by
\(\wedge,\vee,\to\), and bounded \(\forall,\exists\).  Let \(b_A\) be generated
recursively using the Boolean operations and folds above.

\begin{theorem}[Fixed bounded reflection]
\label{thm:bounded-reflection}
For each fixed bounded formula \(A(\vec x)\), the compiler constructs a
primitive-recursive bit \(b_A(\vec x)\) and an ordinary derivation
\[
  \vdash_{\PRAAx}
  \bigl(b_A(\vec x)=1\leftrightarrow A(\vec x)\bigr).
\]
It also constructs a bit-range derivation for \(b_A\).
\end{theorem}

\begin{proof}
External induction on the finite syntax of \(A\).  The atom case is supplied
by the atom interface.  Conjunction and disjunction use the reflected
subformulas and bit truth tables.  Implication uses decidability of the
antecedent bit.  Bounded universal and existential cases use
\cref{lem:all-reflect,lem:any-reflect}.  The recursion generates a concrete
\(\mathsf{PR}\) symbol and a finite derivation for the fixed \(A\).
\end{proof}

The theorem is exactly the kind of metalevel schema described in
\cref{sec:specialization}; it is not a universal satisfaction predicate.

\subsection{Positive strategy certificates}

Truth bits are enough for invariants and closed decisions.  A proof with
bounded alternation may need the finite choices witnessing truth.  Define a
positive certificate recursively:
\begin{itemize}
  \item a true atom has a trivial certificate;
  \item a conjunction has certificates for both conjuncts;
  \item a disjunction stores a selected true side and its certificate;
  \item \(\exists i<N\,B(i)\) stores an index \(i<N\) and a certificate for
        \(B(i)\);
  \item \(\forall i<N\,B(i)\) stores one certificate for each
        \(i<N\), encoded by a bounded sequence.
\end{itemize}
All certificates are finite because all bounds are numerical terms evaluated
at the current parameters.  For a fixed formula, certificate encoding and
checking are primitive recursive.

\begin{theorem}[Positive strategy soundness]
\label{thm:strategy-sound}
For each fixed bounded formula \(A\), the compiler constructs a validity bit
\(\operatorname{ok}_A(c,\vec x)\) such that
\[
  \operatorname{ok}_A(c,\vec x)=1
  \quad\Longrightarrow\quad A(\vec x)
\]
is derivable in ordinary \(\iPRA\).  Moreover, from
\(b_A(\vec x)=1\) it can construct a bounded canonical certificate by finite
least search.
\end{theorem}

\begin{proof}
Soundness is proved by external induction on \(A\).  Atom and conjunction
cases are immediate.  For a disjunction, the checked tag selects one branch,
and the inductive theorem proves that branch.  For an existential, the checker
verifies the bound and the subcertificate, then existential introduction uses
the stored index.  For a universal, the checker verifies the bounded sequence
of subcertificates; for an arbitrary \(i<N\), sequence lookup and the
inductive theorem yield \(B(i)\), allowing bounded universal assembly.

For construction, use the reflected truth bits.  At a disjunction select the
left side if its bit is one, otherwise the right.  At an existential use
bounded least search for the first true body bit.  At a universal build the
bounded sequence by primitive recursion, recursively constructing each
subcertificate.  Since the formula is fixed, the nested construction is a
fixed primitive-recursive program.
\end{proof}

\subsection{Negative certificates and two-sided decisions}

For robust reflection of games or proof search, one may define dual negative
certificates:
\begin{itemize}
  \item falsity of a conjunction selects a false conjunct;
  \item falsity of a disjunction contains negative certificates for both;
  \item falsity of a universal stores one counterexample;
  \item falsity of an existential stores negative certificates for every
        bounded index.
\end{itemize}
Because the formula is decidable, exactly one side can be generated.  This
yields a finite two-sided strategy interface, useful when a local proof
branches on whether a bounded condition holds.  It does not decide unbounded
formulas.

\subsection{Alternation without nested ad hoc searches}

Consider
\[
  \forall i<I\;\exists j<J(i)\;\forall k<K(i,j)\;R(i,j,k).
\]
A naive proof may nest separate search routines and then struggle to show that
the chosen witnesses fit one object-level term.  The fixed strategy compiler
instead generates one certificate type:
\begin{itemize}
  \item for every \(i<I\), a chosen \(j<J(i)\);
  \item for every \(k<K(i,j)\), an atomic certificate for \(R(i,j,k)\).
\end{itemize}
One checker verifies the whole finite strategy.  The public theorem can seal
this strategy if only the final conclusion matters, observe selected outputs,
or name the strategy through a bounded code if required.

\subsection{Example used by the artifact}

The reference test compiles the finite statement
\[
  \forall i<5\;\exists j<4\; (i+j\text{ is even}).
\]
Its canonical certificate chooses witnesses \(0,1,0,1,0\).  The Python
checker accepts the certificate and rejects a tampered witness.  This validates
the finite data algorithm only.  A complete Lean integration must define the
parity atom in \(\mathsf{PR}\), prove its reflection theorem, generate the
strategy checker, and construct the ordinary derivation promised by
\cref{thm:strategy-sound}.
